\FloatBarrier
\section{Sequential disclosure with source-bound decision certificates}\label{sec:supdisclosure}
\FloatBarrier
\subsection{Frozen population, predictions and label domain}
The input is the unchanged November-17 label-pilot archive, SHA256 \texttt{\seqsplit{dba2491c8edd716072987b92a1793394d0892c511f9431d87ff46da69684eb88}}. Its 26 arrays are regenerated from the original R6 Mitra\_GPU and LightGBM prediction payloads. Of 156 array fields, 150 match exactly, including all probabilities, canonical row IDs and labels. Six regenerated logarithm-contrast tables differ by at most $8.89\times10^{-16}$. The immutable archived tables remain the certificate targets. Three result files reproduce byte for byte; budget decisions, regret and stopping counts agree, while regenerated budget endpoints differ by at most $1.14\times10^{-13}$. The selected first split is determined from metadata, before inspecting model gaps or costs. There are 20 AUC tasks and six native multiclass log-loss tasks, containing $43{,}007$ task--row labels. This is a subset of the original 33-task classification/regression audit and is not an additional independent benchmark family.

The pool and its rows stay fixed. In the binary AUC tasks, every unknown row can take either class, subject to at least one positive and one negative in the complete pool. Unknown log-loss labels belong to the same registered class set for both candidates. We make no population-frequency, calibration or exchangeability assumption. The verifier asks which winner would follow from complete disclosure of this pool, rather than estimating a population winner with a confidence interval.

For each unknown log-loss row, independent minimization and maximization of the paired contribution give sharp endpoints, since any combination of its permitted labels is legal. For AUC the positive denominator does not affect the sign. The minimum chooses negative rank-contrast coefficients and the maximum positive coefficients, subject to the disclosed labels. If an unconstrained choice has only one class, forcing the least costly admissible remaining class gives the sharp endpoint. A generic LP with $[0,1]$ label variables and the two class-count inequalities independently checks those endpoints; the constraint matrix gives integral extreme points. The resulting optimization uses established linear methods.

\FloatBarrier
\subsection{Acquisition controls and exact stopping}
The initial policy prior is the mean of the two candidates' probabilities. Maximum entropy sorts its entropy; maximum width sorts the range of the paired contribution; fixed expected gain sorts improvement of the initially favored endpoint. The adaptive expected-gain rule updates its estimated comparison using only labels already disclosed, then targets the corresponding endpoint. Random disclosure uses 50 fixed task-specific permutations with seeds 20261117--20261166. All policies share the same completion domain, unit labels and stopping test.

The implementation separates \nolinkurl{EvidenceState}, \nolinkurl{ClassicalPolicy} and \nolinkurl{ReplayOracle}. The state contains fixed predictions and disclosed labels; the oracle alone holds archived future labels. Every selection is made before the next oracle call. Changing only future labels leaves the observed policy prefix unchanged in 104 checks. All $1{,}404$ frozen paths are replayed by a separate incremental verifier, with the same first certificate and final winner. A further independent implementation imports no production policy, bound or stopping function: it uses direct greedy selection and sorted-cardinality AUC extrema, agreeing on $2{,}808$ checks immediately before and at the saved stopping points. It also verifies all six log-loss hindsight minima. The portable runner changes only the archived program's two filesystem roots; its numerical body is unchanged. The minimum and maximum are represented as Python integers after encoding the finite binary floating-point contrast values with a common power-of-two denominator. AUC coefficients are already integers. This avoids accumulation error in stopping comparisons. It certifies the cached contrast table, not exact real logarithms; the R6 metric reconstruction and source binding remain separate checks.

\begin{table}[t]
\centering\small\setstretch{1.02}\setlength{\tabcolsep}{3.5pt}
\caption{Mean fraction of archived labels disclosed to the same strict certificate (percent; task-equal weighting). Policies are classical controls.}
\label{tab:disclosurepolicies}
\begin{tabularx}{\textwidth}{@{}XN{0.16\textwidth}N{0.16\textwidth}N{0.18\textwidth}@{}}
\toprule
\PAIRTableHead{Policy} & \PAIRTableHead{All 26} & \PAIRTableHead{20 AUC} & \PAIRTableHead{6 log loss}\\\midrule
Random (50 seeds) & 95.76 & 95.16 & 97.73\\
Predictive entropy & 92.15 & 90.15 & 98.83\\
Interval width & 77.97 & 74.47 & 89.63\\
Fixed expected gain & 68.67 & 59.87 & 97.99\\
Adaptive expected gain & \textbf{65.44} & \textbf{59.21} & \textbf{86.20}\\
\bottomrule
\end{tabularx}
\end{table}

\begin{table}[!htbp]
\centering\small\setstretch{1.02}\setlength{\tabcolsep}{3.5pt}\renewcommand{\arraystretch}{1.04}
\caption{All eligible frozen pools. R and A denote random and adaptive expected-gain label percentages. Every complete-pool winner is certified by the replay.}
\label{tab:disclosuretasks}
\begin{tabularx}{\textwidth}{@{}XL{0.15\textwidth}N{0.12\textwidth}N{0.12\textwidth}N{0.12\textwidth}@{}}
\toprule
\PAIRTableHead{Task} & \PAIRTableHead{Metric} & \PAIRTableHead{Rows} & \PAIRTableHead{R (\%)} & \PAIRTableHead{A (\%)}\\\midrule
\nolinkurl{anneal} & Log loss & 300 & 99.09 & 97.00\\
\nolinkurl{Bank_Customer_Churn} & AUC & 3334 & 94.58 & 53.57\\
\nolinkurl{blood-transfusion-service-center} & AUC & 250 & 83.38 & 36.40\\
\nolinkurl{churn} & AUC & 1667 & 98.73 & 84.70\\
\nolinkurl{coil2000_insurance_policies} & AUC & 3274 & 99.73 & 82.53\\
\nolinkurl{credit-g} & AUC & 334 & 95.49 & 68.86\\
\nolinkurl{diabetes} & AUC & 256 & 88.57 & 42.58\\
\nolinkurl{E-CommereShippingData} & AUC & 3667 & 99.02 & 71.26\\
\nolinkurl{Fitness_Club} & AUC & 500 & 91.72 & 58.00\\
\nolinkurl{hazelnut-spread-contaminant-detection} & AUC & 800 & 89.11 & 38.88\\
\nolinkurl{heloc} & AUC & 3487 & 91.86 & 58.76\\
\nolinkurl{in_vehicle_coupon_recommendation} & AUC & 4228 & 98.43 & 78.22\\
\nolinkurl{Is-this-a-good-customer} & AUC & 575 & 93.87 & 53.39\\
\nolinkurl{Marketing_Campaign} & AUC & 747 & 96.48 & 49.67\\
\nolinkurl{maternal_health_risk} & Log loss & 338 & 94.24 & 77.22\\
\nolinkurl{NATICUSdroid} & AUC & 2497 & 97.91 & 50.62\\
\nolinkurl{online_shoppers_intention} & AUC & 4110 & 97.78 & 55.91\\
\nolinkurl{polish_companies_bankruptcy} & AUC & 1970 & 97.86 & 45.43\\
\nolinkurl{qsar-biodeg} & AUC & 352 & 92.55 & 49.43\\
\nolinkurl{seismic-bumps} & AUC & 862 & 98.34 & 65.55\\
\nolinkurl{splice} & Log loss & 1064 & 99.81 & 98.87\\
\nolinkurl{students_dropout_and_academic_success} & Log loss & 1475 & 97.81 & 76.00\\
\nolinkurl{taiwanese_bankruptcy_prediction} & AUC & 2273 & 99.86 & 66.48\\
\nolinkurl{website_phishing} & Log loss & 451 & 98.26 & 87.14\\
\nolinkurl{MIC} & Log loss & 567 & 97.15 & 80.95\\
\nolinkurl{jm1} & AUC & 3629 & 98.04 & 73.93\\
\bottomrule
\end{tabularx}
\end{table}


\FloatBarrier
\subsection{Prior sensitivity and hindsight information}
Prior ablations retain the model scores, labels, pool, metric and certificate and change only the policy probability proxy. Across 26 pools, pooled/Mitra/LightGBM/uniform probabilities require $65.44/63.81/73.06/77.97\%$ of labels. The Mitra prior improves 17 pools, ties two and worsens seven relative to the pooled prior. This is post hoc sensitivity; no best prior was prospectively selected. A misleading prior may reduce efficiency, while the same universal completion test continues to protect the fixed-pool decision.

For log loss, an answer-aware minimum is obtained by sorting each row's actual improvement to the endpoint supporting the full-pool winner. A prefix first exceeding the strict margin is feasible; every smaller subset has at most that prefix's total improvement. This is an exact hindsight cardinality minimum on the independent label domain. It averages $78.09\%$ across six tasks versus $86.20\%$ for the adaptive rule. Knowledge of unobserved answers makes this a lower reference, not an executable premeasurement strategy. We do not extend this box-domain formula to the coupled AUC domain.

Task-level sensitivity is summarized by 20,000 whole-task bootstrap resamples (seed 20261117). The central 95\% interval for the mean absolute label-fraction reduction is 24.60--35.88 percentage points; the relative-saving interval is 25.86--38.24\%. These are exploratory resampling summaries of the 26 frozen pools, not coverage guarantees for future tasks.

\FloatBarrier
\subsection{Scope of the VMA control and other active-selection methods}
VMA~\cite{hara2024vma} estimates pairwise test-loss differences using variance-aware acquisition. We retain its official \nolinkurl{q_by_loss_diff} numerical body at Git blob \nolinkurl{b7edc7ab0eb07a61f47df7584ebaae508964757c}, with $\epsilon=10^{-2}$ and \nolinkurl{normalize=False}. The archived compatibility loader removes only unavailable Numba decorators. A paired log-loss tensor provides the same clipped cache loss. Independent exponential clocks implement weighted sampling without replacement; the matched random control uses equal weights and the same clock seeds. Their sampling distribution agrees with successive remaining-weight draws, while the random-number stream differs from the official rejection sampler.

The 300 VMA-function and 300 matched-random paths cover all six log-loss tasks, with the identical completion certificate. The official query-function adapter is re-executed: all 600 paths and both result CSVs reproduce byte for byte. This control concerns its acquisition function under the common strict certificate, rather than the full VMA algorithm or LURE estimator. The original pilot's random log-loss mean ($97.73\%$) and the matched clock-sampler mean ($97.75\%$) are distinct random realizations. MODEL SELECTOR~\cite{okanovic2025selector} and CODA~\cite{kay2025coda} address efficient model choice through their own acquisition and recommendation procedures. Their published performance numbers are not directly comparable to this strict all-completions stopping criterion, and their algorithms are not run here. Label-efficient model choice and minimum sufficient certificates are prior problems, including the selective-prediction treatment of Kuboyama~\cite{kuboyama2026labels}. The present experiment tests their integration with bound knowledge assertions.

\begin{table}[t]
\centering\small\setstretch{1.02}\setlength{\tabcolsep}{3.5pt}
\caption{Six-task log-loss control (percent labels). VMA denotes the published query function under the common strict certificate, with its matched random sampler. Hindsight uses all answers.}
\label{tab:loglosscontrols}
\begin{tabularx}{\textwidth}{@{}XN{0.13\textwidth}N{0.15\textwidth}N{0.15\textwidth}N{0.15\textwidth}@{}}
\toprule
\PAIRTableHead{Task} & \PAIRTableHead{VMA} & \PAIRTableHead{Matched R} & \PAIRTableHead{Adaptive} & \PAIRTableHead{Hindsight}\\\midrule
\nolinkurl{anneal} & 98.99 & 99.14 & 97.00 & 96.33\\
\nolinkurl{maternal_health_risk} & 95.41 & 94.31 & 77.22 & 67.46\\
\nolinkurl{splice} & 99.74 & 99.81 & 98.87 & 92.39\\
\nolinkurl{students_dropout_and_academic_success} & 97.43 & 97.89 & 76.00 & 58.24\\
\nolinkurl{website_phishing} & 98.24 & 98.27 & 87.14 & 79.16\\
\nolinkurl{MIC} & 96.51 & 97.07 & 80.95 & 74.96\\
\bottomrule
\end{tabularx}
\end{table}


\FloatBarrier
\subsection{Semantic verification and queryable budgets}
Each new assertion binds the native task/repeat/fold/sample key, prediction-array hash, label universe, metric, tolerance, prior, action type, evidence identities, cost unit and exact endpoints. The checker reconstructs the policy path from the frozen source and verifies the first stopping point. It rejects 104 rehashed forgeries changing cost units, pool size, class count or completion-domain semantics. It also rejects 26 requests to treat a new wet measurement as exact disclosure of an old label. Recomputing an object hash alone cannot establish those predicates.

The 26 certificates extend the 896-record core graph to 922 assertions. The new export includes explicit RDF budget fields. A SPARQL query returns winner, labels disclosed, pool size, bounds and cost scope; a separate SQLite ledger returns exactly the same 26 rows. The two query implementations consume the same verified numerical witnesses, so their agreement is a representation check. Numerical endpoint independence is supplied by the LP controls. The original four-family 9,665-operation comparison remains the core system-control result.

\FloatBarrier
\section{Crossed target and observation-measure sensitivity}\label{sec:suptargetmeasure}
The original pilot crosses FULL and within-cell ID targets with original occurrence weights and the observed duplicate-join weights. It uses all 29 non-naive models and ten frozen DrEval folds. The support contains 112,066 original occurrences in 109,298 identity keys: 106,696 singleton keys and 2,602 repeated keys. Repeated-key sufficient statistics preserve the native observations; no per-occurrence identity is invented. The 555 fold--cell groups include 15 with nonuniform lookup multiplicity.

This design distinguishes a declared target change from a changed measure and from refitting a measure-dependent target. The native ID target centers within the original empirical measure. A second computation refits the target after the real join. Holding the original target fixed under joined scoring separates these operations. Inverse-multiplicity weighting with consistent target fitting restores the original measure; inverse weighting only at the final scoring stage need not restore a target already refitted under joined weights.

\begin{table}[t]
\centering\small\setstretch{1.02}\setlength{\tabcolsep}{3.5pt}
\caption{Strict pair-order sensitivities in the frozen 29-model, ten-fold crossed analysis. Undefined Pearson scores are omitted; means require ten defined folds.}
\label{tab:crossedmeasure}
\begin{tabularx}{\textwidth}{@{}L{0.25\textwidth}XL{0.13\textwidth}N{0.1\textwidth}N{0.11\textwidth}@{}}
\toprule
\PAIRTableHead{Before} & \PAIRTableHead{After} & \PAIRTableHead{Metric} & \PAIRTableHead{Fold} & \PAIRTableHead{Mean}\\\midrule
FULL native & FULL joined & MSE & 24 & 2\\
FULL native & FULL joined & Pearson & 19 & 0\\
ID native & ID refit joined & MSE & 19 & 1\\
ID native & ID refit joined & Pearson & 20 & 1\\
FULL native & ID native & MSE & 207 & 19\\
FULL native & ID native & Pearson & 135 & 10\\
FULL joined & ID refit joined & MSE & 200 & 18\\
FULL joined & ID refit joined & Pearson & 140 & 11\\
Fixed ID joined & ID refit joined & MSE & 0 & 0\\
Fixed ID joined & ID refit joined & Pearson & 0 & 0\\
ID native & Refit ID restored weights & MSE & 0 & 0\\
ID native & Refit ID restored weights & Pearson & 0 & 0\\
\bottomrule
\end{tabularx}
\end{table}


The four-condition MSE comparison has 24 foldwise and two mean join changes under FULL, versus 19 and one under ID. The corresponding Pearson counts are 19/0 and 20/1. Switching FULL to ID changes 207/19 MSE fold/mean comparisons under native occurrence weights and 200/18 under joined weights; the Pearson counts are 135/10 and 140/11. The refitted-target and fixed-target joined controls have no additional strict pair-order changes in this snapshot, although scores differ. The maximum MSE difference in the inverse-weight-after-refit control is $1.6451\times10^{-4}$. Forty independent least-squares checks agree to $8.88\times10^{-16}$. These are crossed reanalyses of the existing audit, not a second discovered historical software defect or a new projector theorem. They supply a practical transport check for the main-text observation-measure preservation proposition.

\FloatBarrier
\section{PAX6: qualifying source identity and evidence actions}\label{sec:suppax6}
\FloatBarrier
\subsection{Published scores and reference-bound predictions}
The PAX6 yeast-one-hybrid study~\cite{mcdonnell2024pax6} supplies real barcode-to-variant identities and pooled growth counts. The final source-qualification archive is retained unchanged (SHA256 \texttt{\seqsplit{e42c55c2cb9bf9cd3159191b63e8d5b34d7da847e1e5d5b6a69a4815e87dc41c}}); all 61 internal entries are checked and its original qualification program is executed offline.

Four MaveDB score sets contain 30,344 records, including multi-substitution variants. Against the 422-residue assay reference, every exported protein substitution has the correct reference residue. Across the four Figure-5 source panels, 1,062 matched variant scores agree to at most $5.00\times10^{-10}$; 661 unmatched figure records are synonymous controls. This is a published-source comparison, not a count-to-fitness producer replay.

\begin{table}[t]
\centering\small\setstretch{1.02}\setlength{\tabcolsep}{3.5pt}
\caption{PAX6 published score-set qualification against the fixed 422-residue assay reference. These are records, not independently priced experimental actions.}
\label{tab:pax6qualified}
\begin{tabularx}{\textwidth}{@{}XN{0.12\textwidth}N{0.17\textwidth}N{0.14\textwidth}N{0.13\textwidth}@{}}
\toprule
\PAIRTableHead{Score set} & \PAIRTableHead{All} & \PAIRTableHead{Single missense} & \PAIRTableHead{Single stop} & \PAIRTableHead{Multiple}\\\midrule
pax6\_le9\_geneticin & 7764 & 2709 & 146 & 4909\\
pax6\_blx\_geneticin & 7391 & 2693 & 144 & 4554\\
pax6\_le9\_no\_geneticin & 7819 & 2710 & 147 & 4962\\
pax6\_blx\_no\_geneticin & 7370 & 2687 & 145 & 4538\\
\bottomrule
\end{tabularx}
\end{table}


The EVE file has 8,440 rows over 422 positions and nonmissing predictions for 2,655 of the 2,693 BLX geneticin single-missense variants. The retrieved ESM human MANE file instead describes a 436-residue reference with a 14-residue insertion after residue 47. A second, 461-residue candidate contains no nonmissing ESM scores. Matching residue numbers across these versions would not establish variant identity. These retrieved files therefore do not supply a qualified two-model comparison on the frozen assay. We calculate neither Spearman model superiority nor a model winner, and make no assertion that an aligned ESM resource cannot exist elsewhere.

\FloatBarrier
\subsection{Native samples, shared returns and action semantics}
The consensus table contains 73,144 unique barcodes. Four count tables contain 225,348 records, with 62,298 union and 49,860 intersection barcodes. Reverse-complementing the consensus strand matches every counted barcode. Shared T0 values agree exactly for 52,778 BLX and 57,993 LE9 records. These are source-representation and shared-input checks, rather than independent biological replications.

The 15 GEO source samples comprise one PacBio barcode-phasing sample and seven functional time/condition samples per strain. A sample identifier is not a marginally purchasable action. PacBio phasing and indexed NextSeq sequencing return many records together; shared T0 and preparation may span multiple outputs. The action ledger therefore distinguishes digital recovery of a historical table, a new pooled sequencing experiment, and computational reprocessing. Historical movie duration and culture time points do not establish marginal prices or repeat turnaround.

For exact digital recovery, the observed old field intersects an existing completion set. For new counts $C_{\rm new}$, there is no supplied equality $C_{\rm new}=C_{\rm old}$. Substituting them into the original score can change the outcome being evaluated. A new-measurement study needs a validated observation model, available material, a legal joint identity/count domain and an action-specific cost ledger. The supplied \nolinkurl{variant_call_support} is a call-support count, not a calibrated identity probability; the published T0 threshold is also retained. Under these source records the PAX6 example stops at action qualification. It reports neither mixed-evidence optimization nor experimental-cost savings. This explicit typed boundary prevents an archived-record budget from being reused as an unsupported experimental budget.
